% Appendix A — Wheat Quality Thresholds and Classification Categories

\chapter{Wheat Quality Thresholds and Classification Categories}

\label{AppendixA}

\lhead{Appendix A. \emph{Grain Classification Thresholds}}

This appendix documents the twelve grain and impurity categories central to standard wheat quality assessment protocols, and the corresponding $W/W\%$ thresholds against which each sample is evaluated. These thresholds were sourced from food analysts during the contextual inquiry phase and form the benchmark against which the Drishti classifier is calibrated.

\vspace{1em}

\begin{table}[h!]
\centering
\caption{Wheat grain classification categories and indicative $W/W\%$ thresholds as applied in procurement protocols. Specific threshold values are proprietary and represented here in generalised form.}
\label{tab:thresholds}
\begin{tabular}{llp{5cm}}
\toprule
\textbf{Category} & \textbf{Type} & \textbf{Notes} \\
\midrule
Pure Wheat          & Desired Output    & No upper threshold; higher $\%$ indicates quality. \\
Broken Wheat        & Affected Grain    & $W/W\%$ must remain below threshold. \\
Kernel Burnt        & Affected Grain    & Scorched pericarp; often ventral-side defect. \\
Infested            & Affected Grain    & Insect-bored or contaminated; ventral-side signature. \\
Black Tip           & Affected Grain    & Discolouration at embryo end; subtle chromatic marker. \\
Potia               & Affected Grain    & Immature or shrivelled kernel. \\
Shrunken            & Affected Grain    & Undersized due to growth conditions. \\
Sprouted            & Affected Grain    & Germination initiated; structural compromise. \\
Stone               & Impurity          & High density; settles during vibration. \\
Earth Balls         & Impurity          & Clumped soil; density-variable. \\
Chaff               & Impurity          & Low density; surface-presenting under vibration. \\
Other Impurities    & Impurity          & Miscellaneous non-wheat material. \\
\bottomrule
\end{tabular}
\end{table}

\vspace{1em}

The twelve-category framework above illustrates the classification complexity that the Drishti system is designed to replicate. The concentration of defects with ventral-side signatures (Kernel Burnt, Infested, Black Tip) directly motivated the development of the electromagnetic solenoid grain dispersion system documented in Chapter~\ref{Chapter5}.
